%% kapa: exam=2 round=21 year=2010 period=3 subject=compensation_law subject_ko=감정평가 및 보상법규 seq=3 number=3 points=30
%% source: Q-Net 감정평가사 기출문제 (cst00309) artlSeq=5024183 file=extracted/1008/5024183/2010년 제21회 감정평가사 2차시험 기출문제/3교시.hwp
%% conversion: kordoc 4.24.51 -> tools/gen2.py (auto); review=proofread-sonnet; flags=none
\begin{longquestion}{【문제 3】}
감정평가업자 P와 건설업자 Q는 평소에 친밀한 관계를 유지하고 있다. P는 Q의 토지를 평가함에 있어 친분관계를 고려하여 Q에게 유리하게 평가하였다. 국토해양부장관은 P의 행위가「부동산 가격공시 및 감정평가에 관한 법률」을 위반하였다고 판단하여 과징금, 벌금 또는 과태료의 부과를 검토하고 있다. (30점) \par
\medskip\noindent (1) 과징금, 벌금, 과태료의 법적 성질을 비교하여 설명하시오. (20점) \par
\medskip\noindent (2) 국토해양부장관은 과징금과 벌금을 중복하여 부과하고자 한다. 중복부과처분의 적법성에 관하여 판단하시오. (10점) \par
\end{longquestion}
